
\section{Peyre's tensorization in finite form}\label{app:peyre}

\paragraph{What this appendix proves and how to read it.}
This appendix proves \cref{thm:peyre-finite}, the case of Peyre's
tensorization theorem \cite[Theorem~3.3.1]{Peyre} that the local-law
estimates use: finitely many random variables with finitely many values. If
each variable of one family and each variable of a disjoint
family have conditional maximal correlation at most
$\eps_{ij}<1$ (the numbers in the matrix $\mathcal E$ below),
given any values of any of the other variables, then the two
families have maximal correlation at most the
operator norm of the matrix of these bounds. The paper applies the theorem
in the pinned influence theorem (\cref{prop:middle-negative}), which passes
from the conditional correlations of pairs of vertices of a strict-square
graph law to those of disjoint sets of vertices, and through it in the
correlation bounds of the construction, among them the row correlations of
the twelve-label middle reference and the conditional influence bound of the
upper range. The conditional decorrelation inequality stated with the
theorem (\cref{cor:cond-decorrelation}) is part (b) of \cref{lem:peyre:basic}
below, and the corollary for independent pairs (\cref{cor:peyre-independent})
is the diagonal case (\cref{rem:peyre:sharp}(i)). Only the first statement of
the theorem is used; the second, the prefix form, needs no additional argument, because the
proof starts from it.

The argument is Peyre's \cite[\S3.3.a]{Peyre}. In the finite setting no
approximation argument is needed. We replace the inductive bookkeeping of
\cite[Lemma~3.3.7]{Peyre} by a monotone potential, and the Perron--Frobenius
step \cite[Lemma~3.8.1]{Peyre} by an explicit Neumann series. We also record,
following \cite[\S3.4.a]{Peyre}, that only conditionings on initial segments
of two fixed orderings are used. \Cref{rem:peyre:source} compares the proof
with the source. The proof supplies the finite theorem used in this paper,
independently of the publication status of the source.

For a real matrix $\mathcal E=(\eps_{ij})_{i\in A,\,j\in B}$,
$\opnorm{\mathcal E}$ is its norm as an operator $\ell^2(B)\to\ell^2(A)$
(Euclidean norms, as in \cref{def:prelim:operators}), i.e.\ its largest
singular value. Permuting rows or columns does not change it. As in Peyre's
paper, all letters in this appendix are local to it; in particular $A$ and
$B$ are index sets, and $N$ and $M$ are their cardinalities.

\begin{theorem*}[\Cref{thm:peyre-finite}, restated: tensorization, including the prefix form]
Let $A,B$ be disjoint nonempty finite sets, and let $(X_v)_{v\in A\cup B}$
be random variables on a common probability space, each taking finitely many
values. Let $\eps_{ij}\ge0$ for $i\in A$, $j\in B$, and
$\mathcal E=(\eps_{ij})$. Suppose that for all $i\in A$, $j\in B$, every
$S\subseteq(A\cup B)\setminus\{i,j\}$ and every value $s$ of
$X_S=(X_v)_{v\in S}$ with $\Pr(X_S=s)>0$,
\begin{equation}\tag{\ref*{eq:peyre:hyp}}
  \rho(X_i;X_j\mid X_S=s)\le\eps_{ij}.
\end{equation}
Then
\begin{equation}\tag{\ref*{eq:peyre:concl}}
  \rho\big((X_a)_{a\in A};(X_b)_{b\in B}\big)\le\min\{1,\opnorm{\mathcal E}\}.
\end{equation}
More precisely, \eqref{eq:peyre:concl} holds as soon as there are
enumerations $A=\{a_1,\dots,a_N\}$ and $B=\{b_1,\dots,b_M\}$ such that
\eqref{eq:peyre:hyp} holds for $i=a_k$, $j=b_l$ and the single set
$S=\{a_1,\dots,a_{k-1}\}\cup\{b_1,\dots,b_{l-1}\}$, for all $1\le k\le N$
and $1\le l\le M$.
\end{theorem*}

The second statement has a weaker hypothesis, because the sets $S$ it uses
are among those allowed in the first: they avoid $a_k$ and $b_l$.

\subsection{Conventions}\label{sec:peyre:conv}

\emph{Reduction to a finite space.} Every quantity in the theorem is a
function of the joint law of the finite tuple $(X_v)_{v\in A\cup B}$.
Indeed, square-integrable functions of the $X_v$ are functions on the finite
set of values of the tuple, and conditional laws are taken only given events
of positive probability. We may therefore replace the probability space by
the finite set $\Omega$ of values of the tuple that have positive
probability, with the induced law, and each $X_v$ by the corresponding
coordinate function on $\Omega$. \emph{Throughout this appendix,
$(\Omega,\Pr)$ is a finite probability space with $\Pr(\{\omega\})>0$ for
every $\omega\in\Omega$}, and all functions are real valued. Then
$L^2(\Pr)=\R^\Omega$, with $\ip fg=\E[fg]$ and $\norm f=\E[f^2]^{1/2}$.

\emph{Spaces of functions of variables.} For a random variable $Z$ on
$\Omega$, with values in any set, $L^2(Z)$ is the space of functions of the
form $\omega\mapsto\varphi(Z(\omega))$. Equivalently, it is the space of
functions that are constant on each level set $\{Z=z\}$. For finitely many
variables $Z_1,\dots,Z_k$ we write $L^2(Z_1,\dots,Z_k)=L^2((Z_1,\dots,Z_k))$;
for the empty list this is the space of constants. We write $\E[f\mid Z]$ for
the orthogonal projection of $f$ onto $L^2(Z)$. Explicitly,
$\E[f\mid Z](\omega)=\E[f\one_{\{Z=Z(\omega)\}}]/\Pr(Z=Z(\omega))$ is the
average of $f$ over the level set of $Z$ containing $\omega$. For the empty
list, $\E[f\mid\,\cdot\,]=\E f$. We use the following standard facts.
\begin{enumerate}[label=(P\arabic*)]
\item\label{P:tower} If $L^2(Z)\subseteq L^2(W)$, then
  $\E[\E[f\mid W]\mid Z]=\E[f\mid Z]$ and
  $\E[\E[f\mid Z]\mid W]=\E[f\mid Z]$. In particular $\E[h\mid W]=h$ for
  $h\in L^2(W)$.
\item\label{P:pyth} $\norm f^2=\norm{\E[f\mid Z]}^2+\norm{f-\E[f\mid Z]}^2$,
  and $f-\E[f\mid Z]\perp L^2(Z)$. Moreover
  $\norm{f-\E[f\mid Z]}\le\norm{f-h}$ for every $h\in L^2(Z)$. In particular
  $\norm{f-\E[f\mid Z]}\le\norm f$, and
  $\norm{f-\E[f\mid Z]}\le\norm{f-\E f}$ since constants lie in $L^2(Z)$.
\item\label{P:restrict} Let $W$ be a random variable and $w$ a value with
  $\Pr(W=w)>0$. Put $\Omega_w=\{W=w\}$, with the conditional law
  $\Pr_w=\Pr(\,\cdot\mid W=w)$ and expectation $\E_w$. Then for every random
  variable $Z$ and every $f$, the restriction of $\E[f\mid Z,W]$ to
  $\Omega_w$ equals $\E_w[f|_{\Omega_w}\mid Z|_{\Omega_w}]$. This holds
  because the level sets of $Z|_{\Omega_w}$ in $\Omega_w$ are exactly the
  level sets of $(Z,W)$ contained in $\Omega_w$, and on each of them both
  functions are the average of $f$ for the same (proportional) weights.
\end{enumerate}

\emph{Maximal correlation.} Maximal and conditional maximal correlations are
those of \cref{def:prelim:maxcorr}. In the present notation, with
$L^2_0(X)=\{f\in L^2(X):\E f=0\}$, $\rho(X;Y)$ is the supremum of
$|\E[fg]|$ over $f\in L^2_0(X)$ and $g\in L^2_0(Y)$ with
$\norm f=\norm g=1$, and $\rho(X;Y)=0$ if $X$ or $Y$ is constant; by
Cauchy--Schwarz, $\rho(X;Y)\le1$. For an event $E$ with $\Pr(E)>0$,
$\rho(X;Y\mid E)$ is the same quantity computed in the finite probability
space $(E,\Pr(\,\cdot\mid E))$ for the restrictions of $X$ and $Y$ to $E$,
and $\rho(X;Y\mid Z=z)=\rho(X;Y\mid\{Z=z\})$.

\emph{Orderings.} For random variables $X_1,\dots,X_N$ and $0\le i\le N$ we
write $X_{\le i}=(X_1,\dots,X_i)$ and $X_{<i}=X_{\le i-1}$. The list
$X_{\le0}$ is empty, so $L^2(X_{\le0})$ consists of the constants and
conditioning on $X_{<1}$ is no conditioning.

\subsection{Elementary facts}

Three elementary consequences of the definition of maximal
correlation are used repeatedly; part~(b) is the conditional
decorrelation inequality applied in the main text.

\begin{lemma}[Projection consequences of maximal correlation]\label{lem:peyre:basic}
Let $X,Y,Z$ be random variables on $\Omega$.
\begin{enumerate}[label=(\alph*)]
\item If $f\in L^2_0(X)$ and $g\in L^2(Y)$, then
  $|\E[fg]|\le\rho(X;Y)\,\norm f\,\norm{g-\E g}$. In particular
  $\norm{\E[f\mid Y]}\le\rho(X;Y)\norm f$ for $f\in L^2_0(X)$.
\item Let $\eps\ge0$, and assume $\rho(X;Y\mid Z=z)\le\eps$ for every $z$
  with $\Pr(Z=z)>0$. If $f\in L^2(X,Z)$ satisfies $\E[f\mid Z]=0$ and
  $g\in L^2(Y,Z)$, then
  $|\E[fg]|\le\eps\,\norm f\,\norm{g-\E[g\mid Z]}\le\eps\norm f\norm g$.
\item If $\lambda\ge0$ and $\norm{\E[f\mid Y]}^2\le\lambda\norm f^2$ for every
  $f\in L^2_0(X)$, then $\rho(X;Y)\le\sqrt\lambda$.
\end{enumerate}
\end{lemma}

\noindent Part (b), the \emph{conditional decorrelation} inequality, is the
statement of \cref{cor:cond-decorrelation}; the proof below is complete in
itself and does not use that corollary.

\begin{proof}
(a) Since $\E f=0$, $\E[fg]=\E[f(g-\E g)]$, and $g-\E g\in L^2_0(Y)$. If $f=0$
or $g=\E g$ there is nothing to prove. Otherwise $f/\norm f$ and
$(g-\E g)/\norm{g-\E g}$ are admissible in the definition of $\rho(X;Y)$,
which gives the inequality. For the second claim put
$h=\E[f\mid Y]\in L^2(Y)$. Then $\E h=\E f=0$ by \ref{P:tower}, and
$\norm h^2=\E[hf]$ because $f-h\perp L^2(Y)$ by \ref{P:pyth}. By the first
claim with $g=h$, $\norm h^2\le\rho(X;Y)\norm f\norm h$.

(b) Let $z$ be a value with $\Pr(Z=z)>0$, and put $\Omega_z=\{Z=z\}$,
$\Pr_z=\Pr(\,\cdot\mid Z=z)$, with expectation $\E_z$ and norm
$\norm\cdot_z$. Write $f=F(X,Z)$ and $g=G(Y,Z)$ for some functions $F,G$. On
$\Omega_z$ we have $f=F(X,z)$, a function of $X$, and $g=G(Y,z)$, a function
of $Y$. Moreover $\E_z f$ is the value of $\E[f\mid Z]=0$ on $\Omega_z$, and
$\E_z g$ is the value of $\E[g\mid Z]$ on $\Omega_z$. Part (a) in the space
$(\Omega_z,\Pr_z)$ gives
\[
  |\E_z[fg]|\le\eps\,\norm f_z\,\norm{g-\E[g\mid Z]}_z .
\]
Since $\E[u]=\sum_z\Pr(Z=z)\E_z[u]$ for every function $u$, Cauchy--Schwarz
for the sum over $z$ gives
\begin{align*}
  |\E[fg]|&\le\eps\sum_z\Pr(Z=z)\norm f_z\norm{g-\E[g\mid Z]}_z\\
  &\le\eps\Big(\sum_z\Pr(Z=z)\norm f_z^2\Big)^{1/2}
  \Big(\sum_z\Pr(Z=z)\norm{g-\E[g\mid Z]}_z^2\Big)^{1/2},
\end{align*}
and the right side equals $\eps\norm f\norm{g-\E[g\mid Z]}$. Finally
$\norm{g-\E[g\mid Z]}\le\norm g$ by \ref{P:pyth}.

(c) Let $f\in L^2_0(X)$ and $g\in L^2_0(Y)$ have norm $1$. Since
$g\in L^2(Y)$ and $f-\E[f\mid Y]\perp L^2(Y)$,
$|\E[fg]|=|\E[\E[f\mid Y]\,g]|\le\norm{\E[f\mid Y]}\le\sqrt\lambda$.
\end{proof}

\subsection{One variable against a list}

The next lemma is \cite[Lemma~3.3.3]{Peyre}, with its proof
\cite[Claims~3.3.4 and 3.3.5]{Peyre}.

\begin{lemma}[One new variable and martingale differences]\label{lem:peyre:one}
Let $X_1,\dots,X_N,Y$ be random variables on $\Omega$ and
$\eps_1,\dots,\eps_N\ge0$. Assume that for every $1\le i\le N$ and every
value $x$ of $X_{<i}$ with $\Pr(X_{<i}=x)>0$,
\begin{equation}\label{eq:peyre:one-hyp}
  \rho(X_i;Y\mid X_{<i}=x)\le\eps_i
\end{equation}
(for $i=1$ this reads $\rho(X_1;Y)\le\eps_1$). Let $f\in L^2(X_{\le N})$, and
for $1\le i\le N$ put
\[
  f_i=\E[f\mid X_{\le i}]-\E[f\mid X_{<i}],\qquad
  f_i^*=\E[f\mid X_{\le i},Y]-\E[f\mid X_{<i},Y].
\]
These successive differences of conditional expectations are
\emph{martingale differences}. Put $w_i=\norm{f_i}^2$ and
$w_i^*=\norm{f_i^*}^2$. Then for every $1\le i\le N$,
\begin{equation}\label{eq:peyre:one}
  w_i^*\ \ge\ (1-\eps_i^2)\,w_i-2\eps_i\sqrt{w_i}\sum_{k=i+1}^N\eps_k\sqrt{w_k}.
\end{equation}
\end{lemma}

\begin{proof}
For $0\le i\le N$ put $r_i=f-\E[f\mid X_{\le i}]$ and
$r_i^*=f-\E[f\mid X_{\le i},Y]$. Since
$f\in L^2(X_{\le N})\subseteq L^2(X_{\le N},Y)$, \ref{P:tower} gives
$r_N=r_N^*=0$. For $1\le i\le N$,
\[
  r_{i-1}=f_i+r_i,\qquad r_{i-1}^*=f_i^*+r_i^* .
\]
Here $f_i\in L^2(X_{\le i})$ and $r_i\perp L^2(X_{\le i})$ by \ref{P:pyth},
so $\norm{r_{i-1}}^2=w_i+\norm{r_i}^2$. Likewise
$f_i^*\in L^2(X_{\le i},Y)$ and $r_i^*\perp L^2(X_{\le i},Y)$, so
$\norm{r_{i-1}^*}^2=w_i^*+\norm{r_i^*}^2$. Put
$D_i=\norm{r_i}^2-\norm{r_i^*}^2$ for $0\le i\le N$. Subtracting the two
identities,
\begin{equation}\label{eq:peyre:wdiff}
  w_i-w_i^*=D_{i-1}-D_i\qquad(1\le i\le N),\qquad D_N=0 .
\end{equation}

\emph{Step 1: $r_i-r_i^*=\E[r_i\mid X_{\le i},Y]$, and hence
$D_i=\norm{\E[r_i\mid X_{\le i},Y]}^2\ge0$.} Indeed
$\E[f\mid X_{\le i}]\in L^2(X_{\le i})\subseteq L^2(X_{\le i},Y)$, so by
\ref{P:tower}
\[
  r_i-r_i^*=\E[f\mid X_{\le i},Y]-\E[f\mid X_{\le i}]
  =\E\big[f-\E[f\mid X_{\le i}]\,\big|\,X_{\le i},Y\big]
  =\E[r_i\mid X_{\le i},Y].
\]
Hence $r_i^*=r_i-\E[r_i\mid X_{\le i},Y]$, and \ref{P:pyth} gives
$\norm{r_i}^2=\norm{\E[r_i\mid X_{\le i},Y]}^2+\norm{r_i^*}^2$.

\emph{Step 2: $\sqrt{D_{i-1}}\le\eps_i\sqrt{w_i}+\sqrt{D_i}$ for
$1\le i\le N$.} Let $h=\E[r_{i-1}\mid X_{<i},Y]$. By Step~1 (with $i-1$) and
\ref{P:pyth}, $D_{i-1}=\norm h^2=\E[r_{i-1}h]$, and
\[
  \E[r_{i-1}h]=\E[f_ih]+\E[r_ih].
\]
For the first term, $f_i\in L^2(X_i,X_{<i})$, $h\in L^2(Y,X_{<i})$, and
$\E[f_i\mid X_{<i}]=\E[f\mid X_{<i}]-\E[f\mid X_{<i}]=0$ by \ref{P:tower}.
By \eqref{eq:peyre:one-hyp}, \cref{lem:peyre:basic}(b) applies with $X=X_i$,
$Z=X_{<i}$ and $\eps=\eps_i$, and gives
$|\E[f_ih]|\le\eps_i\sqrt{w_i}\,\norm h$. For the second term,
$h\in L^2(X_{<i},Y)\subseteq L^2(X_{\le i},Y)$, so by \ref{P:pyth} and
Step~1
\[
  |\E[r_ih]|=|\E[\E[r_i\mid X_{\le i},Y]\,h]|\le
  \norm{\E[r_i\mid X_{\le i},Y]}\,\norm h=\sqrt{D_i}\,\norm h .
\]
Thus $\norm h^2\le(\eps_i\sqrt{w_i}+\sqrt{D_i})\norm h$, which gives Step~2
(trivially if $h=0$).

\emph{Step 3: $\sqrt{D_i}\le\sum_{k>i}\eps_k\sqrt{w_k}$ for $0\le i\le N$.}
This holds for $i=N$ because $D_N=0$, and Step~2 carries it from $i$ to
$i-1$.

\emph{Step 4.} By \eqref{eq:peyre:wdiff}, Step~2 and Step~3,
\[
  w_i-w_i^*=D_{i-1}-D_i\le\big(\eps_i\sqrt{w_i}+\sqrt{D_i}\big)^2-D_i
  =\eps_i^2w_i+2\eps_i\sqrt{w_i}\sqrt{D_i}
  \le\eps_i^2w_i+2\eps_i\sqrt{w_i}\sum_{k>i}\eps_k\sqrt{w_k},
\]
which is \eqref{eq:peyre:one}.
\end{proof}

In words: the deficit $D_i$ is the squared norm of the part of the residual
$r_i$ that $Y$ explains beyond $X_{\le i}$. It vanishes at $i=N$, its square
root grows by
at most $\eps_i\sqrt{w_i}$ from $i$ to $i-1$ because $f_i$ is conditionally
decorrelated from $Y$ given $X_{<i}$, and its increments are the losses
$w_i-w_i^*$.

\subsection{Positive weights}

The proof weights the $N$ variables of one family by a positive
vector $\alpha$ with $K\alpha\le\lambda\alpha$, where $K$ is built
from the correlation bounds. Such a vector is supplied explicitly
by a Neumann series, with no appeal to Perron--Frobenius theory.

\begin{lemma}[Positive weights from a Neumann series]\label{lem:peyre:weights}
Let $K$ be an $N\times N$ matrix with nonnegative entries and
$\lambda>\opnorm K$. Then the series
$\alpha=\sum_{m\ge0}\lambda^{-m-1}K^m\one$, where
$\one=(1,\dots,1)^{\mathsf T}$, converges. Its entries satisfy
$\alpha_i\ge\lambda^{-1}>0$, and $K\alpha\le\lambda\alpha$ entrywise.
\end{lemma}

\begin{proof}
The terms satisfy
$\norm{\lambda^{-m-1}K^m\one}\le\lambda^{-1}(\opnorm K/\lambda)^m\sqrt N$, so
the series converges. All terms have nonnegative entries, and the term $m=0$
is $\lambda^{-1}\one$, so $\alpha_i\ge\lambda^{-1}$. Finally
$K\alpha=\sum_{m\ge0}\lambda^{-m-1}K^{m+1}\one=\lambda(\alpha-\lambda^{-1}\one)
=\lambda\alpha-\one\le\lambda\alpha$.
\end{proof}

\begin{lemma}[A monotone potential for nonnegative losses]\label{lem:peyre:potential}
Let $w^0,\ldots,w^M\in[0,\infty)^N$ and let $T_1,\ldots,T_M$ be matrices
with nonnegative entries. Suppose $w^l\ge w^{l-1}-T_lw^{l-1}$ entrywise.
Write $c_k^{(l)}=\sum_i(T_l)_{ik}$ for the column sums. If
$\sum_lc_k^{(l)}\le\lambda<1$ for every $k$, then
\[
 \sum_kw_k^M\ge(1-\lambda)\sum_kw_k^0.
\]
\end{lemma}

\begin{proof}
For $0\le l\le M$, put $C_l=\sum_{m=l+1}^M c^{(m)}$ and
$\Phi_l=\sum_k(1-(C_l)_k)w_k^l$. Then $C_M=0$, $C_{l-1}=C_l+c^{(l)}$, and
$0\le C_l\le\lambda<1$ entrywise. The coefficients in $\Phi_l$ are positive,
so
\begin{align*}
 \Phi_l
 &\ge\sum_k(1-(C_l)_k)(w^{l-1}-T_lw^{l-1})_k\\
 &=\Phi_{l-1}+\sum_k(C_l)_k(T_lw^{l-1})_k
 \ge\Phi_{l-1}.
\end{align*}
Consequently $\sum_kw_k^M=\Phi_M\ge\Phi_0\ge(1-\lambda)\sum_kw_k^0$.
\end{proof}

\noindent The equality in the display uses
$\sum_k(T_lw^{l-1})_k=\sum_kc_k^{(l)}w_k^{l-1}$ and $C_{l-1}=C_l+c^{(l)}$.
No coordinate of $w^l$ is assumed to be monotone in $l$; the potential
$\Phi_l$, which weights each coordinate by one minus the column sums still to
come, is.

\subsection{Proof of the tensorization theorem}

The proof fixes the two enumerations, conditions successively on
initial segments, and shows that a weighted energy cannot lose
more than the factor allowed by $\opnorm{\mathcal E}$.

\begin{proof}[Proof of \cref{thm:peyre-finite}]
\emph{Reduction.} The first hypothesis implies the second one for any
enumerations of $A$ and $B$, as noted after the statement. So fix
enumerations $A=\{a_1,\dots,a_N\}$ and $B=\{b_1,\dots,b_M\}$ for which the
second hypothesis holds. Write $X_k:=X_{a_k}$, $Y_l:=X_{b_l}$,
$\eps_{kl}:=\eps_{a_kb_l}$, $\mathbf X=X_{\le N}$ and $\mathbf Y=Y_{\le M}$.
Since $\{X_S=s\}=\{X_{<k}=x,\ Y_{<l}=y\}$ for
$S=\{a_1,\dots,a_{k-1}\}\cup\{b_1,\dots,b_{l-1}\}$ and $s=(x,y)$, the
hypothesis reads: for $1\le k\le N$, $1\le l\le M$, and every $(x,y)$ with
$\Pr(X_{<k}=x,Y_{<l}=y)>0$,
\begin{equation}\label{eq:peyre:H}
  \rho(X_k;Y_l\mid X_{<k}=x,\ Y_{<l}=y)\le\eps_{kl}.
\end{equation}
We must show $\rho(\mathbf X;\mathbf Y)\le\min\{1,\opnorm{\mathcal E}\}$,
where now $\mathcal E=(\eps_{kl})$ is an $N\times M$ matrix. As explained in
\cref{sec:peyre:conv}, we may assume that $\Omega$ is finite and that all
its points have positive probability. Since $\rho\le1$, we may assume
$\opnorm{\mathcal E}<1$. Fix $\lambda$ with $\opnorm{\mathcal E}^2<\lambda<1$.
The matrix $K=\mathcal E\mathcal E^{\mathsf T}$ has nonnegative entries and
$\opnorm K=\opnorm{\mathcal E}^2<\lambda$. By \cref{lem:peyre:weights} there
is $\alpha\in(0,\infty)^N$ with
\begin{equation}\label{eq:peyre:alpha}
  \sum_{l=1}^M\eps_{kl}\sum_{i=1}^N\eps_{il}\alpha_i
  =(\mathcal E\mathcal E^{\mathsf T}\alpha)_k\le\lambda\alpha_k
  \qquad(1\le k\le N).
\end{equation}
Fix $f\in L^2_0(\mathbf X)$. By \cref{lem:peyre:basic}(c) it suffices to
prove
\begin{equation}\label{eq:peyre:goal}
  \norm{\E[f\mid\mathbf Y]}^2\le\lambda\norm f^2 ;
\end{equation}
letting $\lambda$ decrease to $\opnorm{\mathcal E}^2$ then gives
$\rho(\mathbf X;\mathbf Y)\le\opnorm{\mathcal E}$.

\emph{Step 1: a double martingale decomposition.} For $0\le l\le M$ and
$1\le i\le N$ put
\[
  f_i^l=\E[f\mid X_{\le i},Y_{\le l}]-\E[f\mid X_{<i},Y_{\le l}],\qquad
  w_i^l=\norm{f_i^l}^2 ,
\]
and let $w^l=(w_1^l,\dots,w_N^l)^{\mathsf T}\in[0,\infty)^N$. For fixed $l$
the sum over $i$ telescopes. Since $f\in L^2(X_{\le N},Y_{\le l})$,
\ref{P:tower} gives $\sum_{i=1}^Nf_i^l=f-\E[f\mid Y_{\le l}]$. The summands
are pairwise orthogonal: for $i<k$,
$f_i^l\in L^2(X_{\le i},Y_{\le l})\subseteq L^2(X_{<k},Y_{\le l})$, while
$f_k^l\perp L^2(X_{<k},Y_{\le l})$. The latter holds because $f_k^l$ is the
difference of $f-\E[f\mid X_{<k},Y_{\le l}]$ and
$f-\E[f\mid X_{\le k},Y_{\le l}]$, and both are orthogonal to
$L^2(X_{<k},Y_{\le l})$ by \ref{P:pyth}. Hence
\begin{equation}\label{eq:peyre:sum}
  \sum_{i=1}^Nw_i^l=\norm{f-\E[f\mid Y_{\le l}]}^2\qquad(0\le l\le M).
\end{equation}
For $l=0$ this is $\norm{f-\E f}^2=\norm f^2$. For $l=M$ it is
$\norm{f-\E[f\mid\mathbf Y]}^2=\norm f^2-\norm{\E[f\mid\mathbf Y]}^2$, by
\ref{P:pyth}.

\emph{Step 2: the recursive inequality.} We claim that for $1\le l\le M$ and
$1\le i\le N$,
\begin{equation}\label{eq:peyre:rec}
  w_i^l\ \ge\ (1-\eps_{il}^2)\,w_i^{l-1}
  -\sum_{k=i+1}^N\eps_{il}\eps_{kl}\Big(\frac{\alpha_k}{\alpha_i}\,w_i^{l-1}
  +\frac{\alpha_i}{\alpha_k}\,w_k^{l-1}\Big).
\end{equation}
Fix $l$, and a value $y$ of $Y_{<l}$ with $\Pr(Y_{<l}=y)>0$ (for $l=1$ there
is only the empty value, and $\Omega_y=\Omega$). Put $\Omega_y=\{Y_{<l}=y\}$
and $\Pr_y=\Pr(\,\cdot\mid Y_{<l}=y)$, with expectation $\E_y$. We apply
\cref{lem:peyre:one} in the finite probability space $(\Omega_y,\Pr_y)$, all
of whose points have positive mass. We use the restrictions of
$X_1,\dots,X_N$ and of $Y:=Y_l$, the numbers $\eps_{1l},\dots,\eps_{Nl}$, and
the restriction $\varphi=f|_{\Omega_y}$. Since $f$ is a function of
$\mathbf X$, $\varphi\in L^2(X_{\le N})$ on $\Omega_y$. Hypothesis
\eqref{eq:peyre:one-hyp} holds: if $\Pr_y(X_{<i}=x)>0$, then
$\Pr_y(\,\cdot\mid X_{<i}=x)=\Pr(\,\cdot\mid X_{<i}=x,\,Y_{<l}=y)$, and
\eqref{eq:peyre:H} applies. By \ref{P:restrict} with $W=Y_{<l}$, for every
$0\le i\le N$ the restrictions to $\Omega_y$ satisfy
\[
  \E_y[\varphi\mid X_{\le i}]=\E[f\mid X_{\le i},Y_{<l}]\big|_{\Omega_y},\qquad
  \E_y[\varphi\mid X_{\le i},Y_l]=\E[f\mid X_{\le i},Y_{\le l}]\big|_{\Omega_y}.
\]
Therefore the functions $\varphi_i$, $\varphi_i^*$ of \cref{lem:peyre:one}
are $f_i^{l-1}|_{\Omega_y}$ and $f_i^l|_{\Omega_y}$. Writing
$v_i(y)=\E_y[(f_i^{l-1})^2]$ and $v_i^*(y)=\E_y[(f_i^l)^2]$,
\eqref{eq:peyre:one} gives
\begin{align*}
  v_i^*(y)&\ge(1-\eps_{il}^2)v_i(y)
  -\sum_{k>i}\eps_{il}\eps_{kl}\,2\sqrt{v_i(y)v_k(y)}\\
  &\ge(1-\eps_{il}^2)v_i(y)-\sum_{k>i}\eps_{il}\eps_{kl}
  \Big(\frac{\alpha_k}{\alpha_i}v_i(y)+\frac{\alpha_i}{\alpha_k}v_k(y)\Big).
\end{align*}
The second inequality is $2\sqrt{st}\le cs+t/c$ for $s,t\ge0$, $c>0$, with
$c=\alpha_k/\alpha_i$. Multiply by $\Pr(Y_{<l}=y)$ and sum over $y$. Since
$\sum_y\Pr(Y_{<l}=y)\E_y[u^2]=\E[u^2]$ for every $u$, this gives
\eqref{eq:peyre:rec}.

\emph{Step 3: matrix form.} Let $T_l$ be the $N\times N$ matrix with entries
\[
  (T_l)_{ii}=\eps_{il}^2+\sum_{k>i}\eps_{il}\eps_{kl}\frac{\alpha_k}{\alpha_i},
  \qquad
  (T_l)_{ik}=\eps_{il}\eps_{kl}\frac{\alpha_i}{\alpha_k}\ (k>i),\qquad
  (T_l)_{ik}=0\ (k<i).
\]
Its entries are nonnegative, and \eqref{eq:peyre:rec} says
$w^l\ge w^{l-1}-T_lw^{l-1}$ entrywise. The column sums of $T_l$ are
\[
  c_k^{(l)}:=\sum_{i=1}^N(T_l)_{ik}
  =\eps_{kl}^2+\sum_{i>k}\eps_{kl}\eps_{il}\frac{\alpha_i}{\alpha_k}
  +\sum_{i<k}\eps_{il}\eps_{kl}\frac{\alpha_i}{\alpha_k}
  =\frac{\eps_{kl}}{\alpha_k}\sum_{i=1}^N\eps_{il}\alpha_i\ \ge0,
\]
so that $\sum_{i}(T_lu)_i=\sum_kc_k^{(l)}u_k$ for every $u\in\R^N$. By
\eqref{eq:peyre:alpha}, $\sum_{l=1}^Mc_k^{(l)}\le\lambda$ for every $k$.

\emph{Step 4: apply the potential estimate.} The vectors and matrices from
Steps~1--3 satisfy the hypotheses of \cref{lem:peyre:potential}. Therefore,
by \eqref{eq:peyre:sum},
\[
 \norm f^2-\norm{\E[f\mid\mathbf Y]}^2
 =\sum_kw_k^M\ge(1-\lambda)\sum_kw_k^0
 =(1-\lambda)\norm f^2.
\]
This is \eqref{eq:peyre:goal}. Letting $\lambda$ decrease to
$\opnorm{\mathcal E}^2$ proves the theorem.
\end{proof}

The role of the weights is visible in Step~3: the symmetrization of the
cross terms by $2\sqrt{st}\le cs+t/c$ is chosen so that each column sum of
$T_l$ becomes $(\eps_{kl}/\alpha_k)\sum_i\eps_{il}\alpha_i$, and the
Perron-type inequality $\mathcal E\mathcal E^{\mathsf T}\alpha\le\lambda\alpha$
bounds the sum of these column sums over $l$ by $\lambda$, uniformly in $k$.
Without the weights ($\alpha=\one$) the same argument gives only
$\rho(\mathbf X;\mathbf Y)^2\le\max_k(\mathcal E\mathcal E^{\mathsf
T}\one)_k$, the largest row sum of $\mathcal E\mathcal E^{\mathsf T}$, which
can exceed $\opnorm{\mathcal E}^2$.

\subsection{Remarks}

Two remarks close the appendix: how the proof compares with
Peyre's original argument, and why the conditional form of the
hypothesis cannot be weakened.

\begin{remark}[Comparison with the source]\label{rem:peyre:source}
\Cref{lem:peyre:one} is \cite[Lemma~3.3.3]{Peyre}: its Steps~1 and 2 are
Claims~3.3.4 and 3.3.5 there, and Steps~3--4 are (GO)--(GP). Step~2 of the
proof of the theorem is the passage (HB)--(HF) of \cite{Peyre}, with the
identification (HC)--(HD) of the functions used there. Step~4 replaces the
induction on $M$ in the proof of \cite[Lemma~3.3.7]{Peyre}, and
\cref{lem:peyre:weights} replaces \cite[Lemma~3.8.1]{Peyre}, of which only
the easy inequality ``$\le$'' is needed. That only the conditionings in
\eqref{eq:peyre:H} are used is observed in \cite[\S3.4.a]{Peyre}.
\end{remark}

\begin{remark}[Sharpness and the role of conditioning]\label{rem:peyre:sharp}
(i) The bound can be attained. If the pairs $(X_{a_k},X_{b_k})$,
$1\le k\le N=M$, are independent of each other, then every conditional
correlation in \eqref{eq:peyre:hyp} with $i=a_k$, $j=b_l$ equals
$\rho(X_{a_k};X_{b_k})$ if $k=l$ and $0$ otherwise. So the best choice of
$\mathcal E$ is diagonal and $\opnorm{\mathcal E}=\max_k\rho(X_{a_k};X_{b_k})$.
This is also the value of $\rho(X_A;X_B)$, by the classical tensorization for
independent pairs \cite[Theorem~1.1.19]{Peyre}, which is
\cref{cor:peyre-independent}; its proof also checks the two facts just used.

(ii) The conditional hypothesis cannot be replaced by bounds on unconditional
pairwise correlations. Let $X_1,X_2$ be independent uniform bits and
$Y=X_1+X_2\bmod2$. Then $\rho(X_1;Y)=\rho(X_2;Y)=0$, but $Y$ is a nonconstant
function of $(X_1,X_2)$, so $\rho((X_1,X_2);Y)=1$. Accordingly,
$\rho(X_2;Y\mid X_1=x)=1$, since given $X_1$ the variable $Y$ determines
$X_2$.
\end{remark}

\paragraph{Consequences and what to read next.}
The proof uses only linear algebra in $L^2$ of a finite probability space: orthogonal
projections and the facts \ref{P:tower}--\ref{P:restrict}, the
Cauchy--Schwarz and arithmetic-geometric mean inequalities and a Neumann
series; from the rest of the paper, the definitions of maximal correlation
(\cref{def:prelim:maxcorr}) and of the matrix norm
(\cref{def:prelim:operators}).
